% Chapter 3 -- InitechText (feature topic).  Sources (Law 1):
% os/apps/itext/itext.c (TEXT.EXE: the window "untitled" at (4,44) to
% (588,476), Geneva 9, text 5 px right of / 4 px below the content corner;
% the File and Edit menus and their keys; menus_sync: Cut, Copy and Clear
% only with a selection, Paste only when the Clipboard holds text; Undo,
% Page Setup..., Print... and Print One Copy dimmed; key(): arrows,
% Backspace, printable characters, Return a line end, Ctrl chords to the
% menus only; mouseDown: click, Shift-click, drag selection; Close == Quit;
% guarded(): the Save changes alert; save / save_as: Standard File, CR line
% ends; doc_load: CR LF and LF become CR; say(): the three notices; the
% caret does not blink; the window does not re-wrap when resized);
% os/flair/textedit.h (FLAIR_TE_TEXT_MAX 4096 characters; word wrap);
% os/flair/scrap.c (the Clipboard).  Model: the Mac OS 8.1 SimpleText it
% follows (system7-decomp specs/sys8/simpletext-stdfile.md).  The menu
% entries are registered against TEXT's own bar (flair_app_menus.c).
% GATES: test-flair-itext test-itext-route test-flair-itext-scroll test-stdfile test-scrap test-pubs-flair test-pubs-flair-screens
\Gates{test-flair-itext, test-itext-route, test-flair-itext-scroll,
  test-stdfile, test-scrap, test-pubs-flair, test-pubs-flair-screens}
\chapter{InitechText}
\index{InitechText}

InitechText is a simple text editor.  Use it for memos, notes and lists:
any document that is words and lines, without type styles or pictures.
InitechText saves documents as plain text files, which InitechDOS programs
and the other Initech applications can read.

\section{Opening InitechText}
\Operation{TENANT-LOAD}
\begin{TaskSteps}
\Step Open the \file{APPS} folder and double-click \file{TEXT.EXE}.
\Result{InitechText opens with a new, empty document in a window titled
\texttt{untitled} that fills most of the screen.  The lower menu bar shows
its menus: the Apple menu, \menuname{File} and \menuname{Edit}.}
\Step Type your text.
\Result{What you type appears at the insertion point.  When a line is full,
the next word goes to the start of the next line by itself; press
\key{Enter} only to end a paragraph or to leave a blank line.}
\end{TaskSteps}

\FlairScreen[label=fig:text]{a_text}{A memo in InitechText}

\section{Moving and editing}
\index{InitechText!keys}
\begin{center}\small
\begin{tabular}{@{}p{38mm}p{61mm}@{}}
\textbf{Press} & \textbf{To}\\\hline
letters, digits, punctuation, \key{Space} & insert the character\\
\key{Enter} & end the line\\
\key{Backspace} & erase the character before the insertion point, or the
selected text\\
\key{Left Arrow}, \key{Right Arrow} & move one character\\
\key{Up Arrow}, \key{Down Arrow} & move one line\\
\end{tabular}
\end{center}
There is no key that erases the character after the insertion point, and
\key{Home} and \key{End} are not used.  The backslash key is reserved by
FLAIR (see the \doctitle{FLAIR User's Guide}, Chapter~12), so a backslash
cannot be typed.  The insertion point is a steady vertical line; it does
not blink.

\section{Selecting text}
\index{InitechText!selecting}
To change or move text, first select it.  Selected text is shown white on
black.

\SideHead{Click}Click to move the insertion point to the place you click.
\SideHead{Drag}Drag across text to select it.
\SideHead{Shift-click}Hold down \key{Shift} and click to select everything
between the insertion point and the place you click.
\SideHead{Select All}Choose \menuitem{Select All} from the \menuname{Edit}
menu, or press \key{Ctrl-A}, to select the whole document.

\FlairScreen[label=fig:textsel]{a_textsel}{A word selected}

Typing replaces the selected text with what you type; \key{Backspace}
erases it.

\section{Cutting, copying and pasting}
\index{Clipboard}\index{InitechText!Cut, Copy, Paste}
\begin{TaskSteps}
\Step Select the text you want to copy.
\Step Choose \menuitem{Copy} from the \menuname{Edit} menu, or press
\key{Ctrl-C}.
\Result{A copy of the text is put on the Clipboard.  \menuitem{Cut}
(\key{Ctrl-X}) does the same and removes the text from the document.}
\end{TaskSteps}

\FlairScreen[scale=4, label=fig:textedit1]{a_textedit1}{The Edit menu with
text selected}

\begin{TaskSteps}[resume]
\Step Click where you want the text to go.
\Step Choose \menuitem{Paste} from the \menuname{Edit} menu, or press
\key{Ctrl-V}.
\Result{The text on the Clipboard is put at the insertion point, or in
place of the selected text.}
\end{TaskSteps}

\FlairScreen[label=fig:textpaste]{a_textpaste}{The copied word, pasted on
a new line}
\Operation{TENANT-SCRAP}

The text stays on the Clipboard until you cut or copy something else, even
after you quit InitechText, so you can paste it into another document
later.  The Clipboard is emptied when the computer is restarted.

\section{Saving a document}
\index{InitechText!saving}
\Operation{TENANT-SF}
Choose \menuitem{Save} from the \menuname{File} menu, or press
\key{Ctrl-S}.  The first time you save a document, InitechText asks for a
name in the Save dialog box (the \doctitle{FLAIR User's Guide}, Chapter~5,
describes it).

\FlairScreen[label=fig:textsave]{a_textsave}{Saving a new document}

Type the name, for example \type{memo.txt}, and press \key{Enter}.  The
window takes the name of the file.

\FlairScreen[label=fig:textnamed]{a_textnamed}{The document saved as
MEMO.TXT}

After that, \menuitem{Save} saves the document under its name without
asking.  \menuitem{Save As...} always asks for a name, so you can keep a
copy under another name; the window then takes the new name.

InitechText saves a document as a text file, with a carriage return at the
end of each paragraph.  When it opens a file whose lines end with a
carriage return and line feed, or a line feed alone, as files written by
InitechDOS programs do, it reads each of them as one line end, and saves
the file with carriage returns.

\section{Opening a document}
\index{InitechText!opening a document}
Choose \menuitem{Open...} from the \menuname{File} menu, or press
\key{Ctrl-O}, select the document in the Open dialog box, and click
\texttt{Open}.

\FlairScreen[label=fig:textopen]{a_textopen}{Opening a document}

InitechText holds one document at a time.  If the document in its window
has unsaved changes, it first asks whether to save them.  A document can
hold up to 4096 characters.  InitechText does not open a longer file, and
does not take typing or pasting that would make a document longer.

\section{Closing and quitting}
\index{InitechText!quitting}
\Operation{TENANT-ALERT}
InitechText has one window; closing it quits InitechText.  Choose
\menuitem{Close} or \menuitem{Quit} from the \menuname{File} menu, press
\key{Ctrl-W} or \key{Ctrl-Q}, or click the close box.  If the document has
changes you have not saved, InitechText asks what to do with them.

\FlairScreen[label=fig:textask]{a_textask}{InitechText asks before quitting}

Click \texttt{Save} to save the document first, \texttt{Don't Save} to quit
without saving, or \texttt{Cancel} to go back to the document.

\section{The File menu}
\FlairScreen[scale=4, label=fig:textfile]{a_textfile}{InitechText's File
menu}
\begin{MenuReference}{TEXT}{File}
\MenuItemEntry{New}{N}{on}{Clears the window for a new, untitled document,
after asking about unsaved changes.}
\MenuItemEntry{Open...}{O}{on}{Opens a document with the Open dialog box,
after asking about unsaved changes.}
\MenuItemEntry{Close}{W}{on}{Closes the window and quits InitechText, after
asking about unsaved changes.}
\MenuItemEntry{Save}{S}{on}{Saves the document under its name; for a new
document, asks for a name first.}
\MenuItemEntry{Save As...}{}{on}{Saves the document under a name you give
in the Save dialog box.}
\MenuItemEntry{Page Setup...}{}{off}{Dimmed.  InitechText does not print.}
\MenuItemEntry{Print...}{P}{off}{Dimmed.}
\MenuItemEntry{Print One Copy}{}{off}{Dimmed.}
\MenuItemEntry{Quit}{Q}{on}{Quits InitechText, after asking about unsaved
changes.}
\end{MenuReference}

\section{The Edit menu}
\FlairScreen[scale=4, label=fig:textedit0]{a_textedit0}{InitechText's Edit
menu with nothing selected}
\begin{MenuReference}{TEXT}{Edit}
\MenuItemEntry{Undo}{Z}{off}{Dimmed.}
\MenuItemEntry{Cut}{X}{off}{Removes the selected text and puts it on the
Clipboard.  Dimmed when no text is selected.}
\MenuItemEntry{Copy}{C}{off}{Puts a copy of the selected text on the
Clipboard.  Dimmed when no text is selected.}
\MenuItemEntry{Paste}{V}{off}{Puts the text on the Clipboard at the
insertion point, or in place of the selection.  Dimmed when the Clipboard
is empty.}
\MenuItemEntry{Clear}{}{off}{Removes the selected text without putting it
on the Clipboard.  Dimmed when no text is selected.}
\MenuItemEntry{Select All}{A}{on}{Selects the whole document.}
\end{MenuReference}

\section{Notes}
\begin{itemize}
\item The window's text is set in a single size of type, Geneva 9 point.
\item The lines are always the same length, whatever the size of the
  window: changing the size of the window does not rearrange the text.
\item Use the scroll bar to see a document longer than the window.
\end{itemize}
